\documentclass{article}
\usepackage[margin=1.15in]{geometry}
\usepackage{xcolor}
\usepackage{parskip}
\usepackage{fancyhdr}
\usepackage{array}
\usepackage{booktabs}
\usepackage{enumitem}
\usepackage{amsmath}
\usepackage{tabularx}

\pagestyle{fancy}
\fancyhf{}
\cfoot{\thepage}
\rhead{Lecture 2}
\lhead{MA 214: Applied Statistics}
\renewcommand{\headrulewidth}{.4mm}

\newcommand{\blankline}[1]{\underline{\hspace{#1}}}

\begin{document}

\noindent\textbf{Name:}\blankline{2.25in}\hfill\textbf{BUID:}\blankline{1.55in}\newline

\begin{center}
 \Large In-Class Worksheet: Statistical Models and Maximum Likelihood
\end{center}

\subsection*{Scenario}
A statistical model gives a probability description of possible data values. Today we practice choosing simple models, writing likelihoods, and connecting maximum likelihood to least squares.

\medskip
\noindent\textbf{Goals:} \textbf{(1)} match data types to simple probability models, \textbf{(2)} compute an MLE for a Bernoulli model (a proportion), \textbf{(3)} compute the MLE of a mean and variance for a normal model.

\subsection*{Part 1: Choose a Model and Name the Parameter}
For each variable, write one possible model and name the parameter(s). There may be more than one reasonable model.

\begin{center}
\renewcommand{\arraystretch}{1.55}
\begin{tabularx}{\textwidth}{p{6.7cm}p{4.2cm}X}
\toprule
\textbf{Variable} & \textbf{One possible model} & \textbf{Parameter(s)} \\
\midrule
Whether a randomly selected student submitted the weekly homework. & & \\[2.1em]
Time until a course webpage loads, measured in seconds. & & \\[2.1em]
Number of emails received by an instructor in one hour. & & \\[2.1em]
Number of quiz objectives passed by one student out of 4 objectives. & & \\
\bottomrule
\end{tabularx}
\end{center}

\subsection*{Part 2: MLE from Start to Finish}
A small survey asks 12 students whether they used public transportation this week.
\[
1,\ 1,\ 0,\ 1,\ 0,\ 1,\ 1,\ 1,\ 0,\ 1,\ 1,\ 0.
\]
Let $X_i = 1$ if student $i$ used public transportation, and $X_i=0$ otherwise.

\begin{enumerate}[label={\bfseries (\Alph*)},itemsep=.75em]
\item Choose a model for $X_i$:
\[
X_i \sim \blankline{2.8in}
\]
\item Identify the parameter and describe it in words. \\[3em]
\item Count the number of successes and failures.
\[
\#\text{ successes}=\blankline{1in}, \qquad \#\text{ failures}=\blankline{1in}.
\]
\item Write the likelihood as a function of $p$.
\[
L(p) = \blankline{4.5in}
\]
\item Write the log-likelihood, ignoring terms that do not depend on $p$.
\[
\ell(p) = \blankline{4.5in}
\]
\item Find the MLE and interpret it in context.
\[
\hat{p}_{\text{MLE}} = \blankline{2.2in}
\]
\vspace{2em}
\end{enumerate}

\subsection*{Part 3: A Normal Model}
Suppose we observe four webpage loading times, in seconds:
\[
4.2,\quad 5.1,\quad 4.7,\quad 6.0.
\]
Assume $X_i \sim N(\mu,\sigma^2)$ independently.

\begin{enumerate}[label={\bfseries (\Alph*)},itemsep=.75em]
\item Compute $\hat\mu_{\text{MLE}}=\bar{x}$.
\[
\bar{x}=\blankline{3in}
\]
\item Compute the MLE of the variance. Remember: the MLE uses denominator $n$, not $n-1$.
\[
\hat\sigma^2_{\text{MLE}} = \frac{1}{n}\sum_{i=1}^n (x_i-\bar{x})^2 = \blankline{3in}
\]
\item In one sentence, what does $\sigma^2$ measure in this model? \\[3em]
\end{enumerate}

\end{document}
